\section*{Aufgabenstellung}
Ermitteln Sie die Informationen zu den Betriebssystemen der beiden Images.
(Wenn verfügbar mit Build Release, Service Pack, Installationsdatum etc.)

\section*{Lösung}

\subsection*{Linux Server Image}
Das Linux Server Image ist verschlüsselt:

\begin{lstlisting}[language=bash]
luis@Alpha-17:~/Programming/GitHub/University/DigitalForensic$ fsstat -o 2048 images/20260328_server.E01
Possible encryption detected (High entropy (7.96))
\end{lstlisting}

Die Infos musste ich mir daher über binwalk besorgen:
\begin{lstlisting}[language=bash]
root@Alpha-17:/home/luis/Programming/GitHub/University/DigitalForensic# ewfinfo images/20260328_server.E01 | grep -iE 'operating|system|date|media'
        Acquisition date:       Mon Mar 23 15:01:32 2026
        System date:            Mon Mar 23 15:01:32 2026
        Operating system used:  Linux
Media information
        Media type:             fixed disk
        Media size:             4.0 GiB (4298928128 bytes)
root@Alpha-17:/home/luis/Programming/GitHub/University/DigitalForensic# strings /mnt/ewf_server/ewf1 | grep -oP 'Linux version \S+'
Linux version 6.12.74+deb13+1-amd64
root@Alpha-17:/home/luis/Programming/GitHub/University/DigitalForensic# strings /mnt/ewf_server/ewf1 | grep -i 'PRETTY_NAME'
PRETTY_NAME="Debian GNU/Linux 13 (trixie)"
\end{lstlisting}

Es handelt sich also um ein Debian GNU/Linux 13 (trixie) mit Linux Kernel Version 6.12.74+deb13+1-amd64.

\subsection*{Windows Image}

Bei dem Windows Image handelt es sich um ein Windows 10 Pro installation.
Folgende Infos konnten aus der Registry extrahiert werden:
\begin{lstlisting}[language=bash]
root@Alpha-17:/home/luis/Programming/GitHub/University/DigitalForensic# regripper -p winver -r /tmp/SOFTWARE
Launching winver v.20200525
winver v.20200525
(Software) Get Windows version & build info

ProductName               Windows 10 Pro
ReleaseID                 2009
BuildLab                  22621.ni_release.220506-1250
BuildLabEx                22621.1.amd64fre.ni_release.220506-1250
CompositionEditionID      Enterprise
RegisteredOrganization
RegisteredOwner           Windows-Benutzer
InstallDate               2026-02-16 13:44:42Z
InstallTime               2026-02-16 13:44:42Z
\end{lstlisting}

Die Reg Datei wurde wie folgt extrahiert:
\begin{lstlisting}[language=bash]
root@Alpha-17:/home/luis/Programming/GitHub/University/DigitalForensic# fls -pr /dev/loop35p3 | grep 'SOFTWARE$'
r/r 101808-128-4:       Windows/System32/config/SOFTWARE
root@Alpha-17:/home/luis/Programming/GitHub/University/DigitalForensic# icat /dev/loop35p3 101808-128-4 > /tmp/SOFTWARE
root@Alpha-17:/home/luis/Programming/GitHub/University/DigitalForensic# file /tmp/SOFTWARE
/tmp/SOFTWARE: MS Windows registry file, NT/2000 or above
\end{lstlisting}
